% ===== 速通讲义 主文档导言区（include-in-header，配 documentclass=ctexrep）=====
% 来自 autopku speed-notes-kit。唯一需按课程改的是下面 fancyhead[L] 的页眉文字。

% --- 代码/等宽字体（需覆盖希腊字母、数学关系符、箭头等，否则代码注释缺字；CJK 仍由 ctex 的 xeCJK 处理）---
% Menlo 是 macOS 字体；Windows 换 Consolas，Linux 换 DejaVu Sans Mono
\setmonofont{Menlo}[Scale=0.9]

% --- 数学 ---
\usepackage{mathtools}
\usepackage{amssymb}
\usepackage{amsthm}

% --- 彩色提示框（callout.lua 生成 tcolorbox 环境）---
\usepackage[most]{tcolorbox}

% --- 页面尺寸 ---
\usepackage{geometry}
\geometry{a4paper, top=2.4cm, bottom=2.2cm, left=2.3cm, right=2.3cm, headheight=15pt}

% --- 关闭所有自动章节编号（编号已手写进标题文本，避免双重编号）---
\setcounter{secnumdepth}{-1}
\setcounter{tocdepth}{1}

% --- 章（每章）起始更紧凑一点的版式 ---
\ctexset{
  chapter = {
    beforeskip = {6pt},
    afterskip  = {18pt},
    format     = {\raggedright\Large\bfseries},
  },
  section = {
    format = {\large\bfseries\color{blue!30!black}},
  },
  subsection = {
    format = {\normalsize\bfseries},
  }
}

% --- 页眉页脚 ---
\usepackage{fancyhdr}
\pagestyle{fancy}
\fancyhf{}
\fancyhead[L]{\small\itshape 课程名 · 速通讲义}% ← 改成本课程名，如「概率统计 · 零基础速通讲义」
\fancyhead[R]{\small\itshape\leftmark}
\fancyfoot[C]{\small\thepage}
\renewcommand{\headrulewidth}{0.4pt}
\renewcommand{\footrulewidth}{0pt}
\fancypagestyle{plain}{%
  \fancyhf{}%
  \fancyfoot[C]{\small\thepage}%
  \renewcommand{\headrulewidth}{0pt}%
}
\renewcommand{\chaptermark}[1]{\markboth{#1}{}}

% --- 段落与行距 ---
\setlength{\parskip}{0.45em plus 0.12em minus 0.05em}
\linespread{1.06}

% --- 代码块用 listings（支持长行自动折行，无需 fvextra；讲义基本无代码，留作兜底）---
\usepackage{listings}
\lstset{
  breaklines=true,
  breakatwhitespace=false,
  basicstyle=\small\ttfamily\upshape,
  commentstyle=\ttfamily\upshape\color{green!38!black},
  stringstyle=\ttfamily\upshape\color{purple!55!black},
  keywordstyle=\ttfamily\upshape\color{blue!60!black},
  identifierstyle=\ttfamily\upshape,
  columns=fullflexible,
  keepspaces=true,
  showstringspaces=false,
  extendedchars=true,
  frame=single,
  framesep=4pt,
  rulecolor=\color{gray!55},
  backgroundcolor=\color{gray!8},
  xleftmargin=4pt,
  xrightmargin=4pt,
  aboveskip=6pt,
  belowskip=6pt,
}

% --- 表格不要太挤 ---
\renewcommand{\arraystretch}{1.18}

% --- 让 boxed 公式不被页面切断时好看一些 ---
\setlength{\fboxsep}{4pt}
